% Section 1.1: 研究背景与意义

近年来，以 GPT、LLaMA、Qwen 等为代表的大语言模型（Large Language Models, LLMs）在自然语言理解、知识问答与代码生成等任务中表现出接近乃至超越传统专用模型的综合能力，推动人工智能从“单点任务求解”迈向“通用语言智能”阶段\cite{brown2020language,radford2019language,zhang2026llm4se,yang2024dlcodegen,cao2025kgllm}。在软件工程领域，这一趋势正在重塑开发者的工作方式：从早期的语法补全、模板生成，发展为基于自然语言描述自动产生可执行代码的智能编程范式。GitHub Copilot、Cursor、通义灵码等工具在工业界的快速普及，已经印证了大语言模型在真实研发场景中的生产力价值，也进一步激发了学术界对“智能代码生成”这一方向的系统性研究兴趣。

在代码生成任务中，算法代码生成是一类具有代表性且具挑战性的子任务\cite{yang2024dlcodegen,li2026constrained}。与业务逻辑代码相比，算法代码对逻辑严密性、边界条件处理与程序可执行正确性的要求更高：模型不仅需要正确理解题目描述中的约束与目标，还需在有限的上下文中完成问题抽象、数据结构选择、算法设计与具体实现的完整链路。这一过程本质上是一种结构化推理行为，对模型的规划能力与组合泛化能力提出了更高要求。

然而，仅依赖通用大模型的零样本生成，在算法题场景中仍面临较为明显的瓶颈：
\begin{itemize}
    \item \textbf{推理过程不稳定}：同一问题在不同采样下可能给出逻辑不一致的解法；
    \item \textbf{复杂算法正确率偏低}：在动态规划、图论、数论等问题上显著下降，错误多集中在边界条件处理、循环不变量维护与隐含约束识别；
    \item \textbf{全量微调不现实}：算力成本与迭代效率限制下，如何在有限资源下实现稳定、有效的任务适配成为落地层面必须解决的问题。
\end{itemize}

在上述背景下，围绕“如何让大语言模型更好地生成算法代码”开展研究具有重要的理论与应用意义。从理论层面，该问题涉及指令微调、参数高效微调、提示工程与结构化推理等多条技术路线的交叉与融合，有助于深化对大模型推理行为与任务适配机制的理解；从应用层面，高质量的算法代码生成能力可以直接服务于智能编程助手、在线评测辅助、程序设计教学等场景，具有明确的工程价值。本文正是在这一背景下，围绕数据构建、参数高效微调与链式思维推理三个关键环节，开展面向算法代码生成的系统性研究与实现工作。